\documentclass[12pt]{article}
\usepackage[utf8]{inputenc}
\usepackage[T1]{fontenc}
\usepackage[margin=1in]{geometry}
\usepackage{setspace}
\usepackage{xcolor}
\usepackage[colorlinks=true, urlcolor=blue!60!black, linkcolor=blue!60!black]{hyperref}
\setlength{\parindent}{2em}
\setlength{\parskip}{1.25em}
\renewcommand{\baselinestretch}{1.0}
\begin{document}
\begin{center}
{\LARGE \textit{Bolling v. Sharpe}, 347 U.S. 497 (1954)\par}
\vspace{0.8em}
{\large Case Brief\par}
\end{center}
\vspace{1em}

\section{Memory Jogger}

Because it would be ``unthinkable'' for the Constitution to forbid racial school segregation by the states while allowing it by the federal government, the Fifth Amendment's Due Process Clause prohibits segregated public schools in Washington, D.C., launching the doctrine of reverse incorporation.

\section{Detailed Case Facts}

In 1947, Washington, D.C. barber and civil rights activist Gardner Bishop co-founded the Consolidated Parents Group to challenge overcrowded, under-resourced segregated schools for Black students in the nation's capital (\href{https://clearinghouse.net/case/11110/}{Civil Rights Litigation Clearinghouse}; \href{https://historicsites.dcpreservation.org/items/show/969}{DC Preservation League}). After the District of Columbia Board of Education denied the group's petition to open the newly constructed John Philip Sousa Junior High School in Anacostia as an integrated facility, Bishop and the Consolidated Parents Group brought eleven Black students---including eleven-year-old Spottswood Thomas Bolling Jr.---to Sousa Junior High on opening day, September 11, 1950, and sought their admission (\href{https://historicsites.dcpreservation.org/items/show/969}{DC Preservation League}; \href{https://www.oyez.org/cases/1940-1955/347us497}{Oyez}). Although the school had empty classrooms, the principal turned the students away solely because of their race in accordance with the Board of Education's segregation policy (\href{https://clearinghouse.net/case/11110/}{Civil Rights Litigation Clearinghouse}; \href{https://supreme.justia.com/cases/federal/us/347/497/}{Justia}).

Following the illness of pioneering civil rights attorney Charles Hamilton Houston, Howard University School of Law professors James M. Nabrit Jr. and George E. C. Hayes took up the case on behalf of Spottswood Bolling and four other Black students against C. Melvin Sharpe, President of the District of Columbia Board of Education, and other school officials (\href{https://historicsites.dcpreservation.org/items/show/969}{DC Preservation League}). Unlike the companion state school-segregation cases consolidated in \textit{Brown v. Board of Education}, 347 U.S. 483 (1954), which attacked both the inequality of segregated facilities and the ``separate but equal'' doctrine of \textit{Plessy v. Ferguson}, 163 U.S. 537 (1896), Nabrit framed \textit{Bolling} as a direct constitutional challenge to racial segregation itself, independent of facility comparisons (\href{https://supreme.justia.com/cases/federal/us/347/497/}{Justia}).

The constitutional obstacle was structural and textual. Because the District of Columbia is a federal enclave governed by Congress rather than a state, the Fourteenth Amendment's Equal Protection Clause---which by its terms binds only ``any State''---did not apply (\href{https://www.courtlistener.com/opinion/105222/bolling-v-sharpe/}{CourtListener}; \href{https://supreme.justia.com/cases/federal/us/347/497/}{Justia}). Instead, petitioners relied primarily on the Due Process Clause of the Fifth Amendment (``nor be deprived of life, liberty, or property, without due process of law''), along with federal civil rights statutes (42 U.S.C. \S\S~1981 and 1983), the Bill of Attainder Clause of Article I, \S~9, and the United Nations Charter (\href{https://clearinghouse.net/case/11110/}{Civil Rights Litigation Clearinghouse}). Historically and politically, the case arrived at the height of the Cold War and alongside the four state cases in \textit{Brown}; upholding segregation in the nation's capital on the very day the Court struck down Jim Crow schools across the South would have created a profound moral and institutional crisis for the federal judiciary.

\section{Procedural History}

On November 9, 1950, petitioners filed a class action for declaratory and injunctive relief in the United States District Court for the District of Columbia (\href{https://clearinghouse.net/case/11110/}{Civil Rights Litigation Clearinghouse}). The District Court dismissed the complaint, relying on the recent D.C. Circuit precedent in \textit{Carr v. Corning}, 182 F.2d 14 (D.C. Cir. 1950), which had upheld the District's dual school system on the view that Congress had long sanctioned segregated schools in Washington, D.C. (\href{https://historicsites.dcpreservation.org/items/show/969}{DC Preservation League}; \href{https://supreme.justia.com/cases/federal/us/347/497/}{Justia}).

Petitioners appealed to the United States Court of Appeals for the District of Columbia Circuit. While the appeal was pending, the U.S. Supreme Court took the extraordinary step of granting a writ of certiorari before judgment in the Court of Appeals, 344 U.S. 873 (1952), citing the importance of the constitutional question and the need to hear the District of Columbia case alongside the four state cases in \textit{Brown} (\href{https://www.courtlistener.com/opinion/105222/bolling-v-sharpe/}{CourtListener}; \href{https://supreme.justia.com/cases/federal/us/347/497/}{Justia}). The case was initially argued on December 10--11, 1952, reargued on December 8--9, 1953, and decided on May 17, 1954---the same day as \textit{Brown v. Board of Education} (\href{https://supreme.justia.com/cases/federal/us/347/497/}{Justia}). The Court held racial segregation in D.C. public schools unconstitutional under the Fifth Amendment and restored the case to the docket for further argument on the form of the remedial decree, which was subsequently resolved in \textit{Brown v. Board of Education} (\textit{Brown II}), 349 U.S. 294 (1955) (\href{https://supreme.justia.com/cases/federal/us/347/497/}{Justia}).

\section{Judicial Votes}

\begin{itemize}
  \item \textbf{Majority (9--0):} Chief Justice Earl Warren (author), joined by Justices Hugo L. Black, Stanley F. Reed, Felix Frankfurter, William O. Douglas, Robert H. Jackson, Harold H. Burton, Tom C. Clark, and Sherman Minton (\href{https://supreme.justia.com/cases/federal/us/347/497/}{Justia}).
  \item \textbf{Concurrences:} None.
  \item \textbf{Dissents:} None.
\end{itemize}

\section{Holding}

\textbf{Judgment:} Racial segregation in the public schools of the District of Columbia violates the Due Process Clause of the Fifth Amendment; the case was restored to the docket for reargument on the appropriate remedial decree (\href{https://www.courtlistener.com/opinion/105222/bolling-v-sharpe/}{CourtListener}; \href{https://supreme.justia.com/cases/federal/us/347/497/}{Justia}).

\textbf{Rule Announced:} Although the Fifth Amendment lacks an explicit Equal Protection Clause, the concepts of due process and equal protection are not mutually exclusive, and federal discrimination may be so unjustifiable as to violate due process. Racial classifications are constitutionally suspect and must be scrutinized with particular care. Because racial segregation in public education is not reasonably related to any proper governmental objective, it imposes an arbitrary deprivation of liberty on Black students in violation of the Fifth Amendment's Due Process Clause (\href{https://supreme.justia.com/cases/federal/us/347/497/}{Justia}).

\textbf{Holding vs.\ Dicta:} The strict holding is narrow: federally mandated racial segregation in District of Columbia public schools violates the Fifth Amendment's Due Process Clause. Chief Justice Warren's broader statements---that equal protection and due process both stem from ``our American ideal of fairness,'' that ``liberty under law extends to the full range of conduct which the individual is free to pursue,'' and that it would be ``unthinkable'' for the Constitution to impose a lesser anti-discrimination duty on the federal government than on the states---functioned as expansive doctrinal rationale (arguably dicta as to non-educational federal classifications) that later decisions elevated into the full-fledged doctrine of ``reverse incorporation,'' culminating in \textit{Adarand Constructors, Inc. v. Pena}, 515 U.S. 200 (1995), which treated Fifth and Fourteenth Amendment equal protection standards as largely congruent (\href{https://virginialawreview.org/wp-content/uploads/2020/12/1879.pdf}{Virginia Law Review}).

\section{Analysis of Opinions}

\subsection{Majority Opinion (Chief Justice Warren)}

Chief Justice Warren's remarkably brief opinion---spanning barely three pages in the United States Reports---moves through four interlocking steps to bridge the textual gap left by the Fourteenth Amendment's limitation to the states (\href{https://supreme.justia.com/cases/federal/us/347/497/}{Justia}).

First, Warren confronts the textual difference directly. He concedes that the Fifth Amendment contains no equal protection clause and clarifies that ``due process of law'' and ``equal protection of the laws'' are not ``always interchangeable phrases,'' with equal protection offering a ``more explicit safeguard of prohibited unfairness'' (\href{https://supreme.justia.com/cases/federal/us/347/497/}{Justia}). Yet he insists the two guarantees are ``not mutually exclusive'' because both stem from ``our American ideal of fairness.'' Citing \textit{Detroit Bank v. United States}, 317 U.S. 329 (1943), \textit{Currin v. Wallace}, 306 U.S. 1 (1939), and \textit{Steward Machine Co. v. Davis}, 301 U.S. 548 (1937), he states the doctrinal bridge: ``discrimination may be so unjustifiable as to be violative of due process'' (\href{https://supreme.justia.com/cases/federal/us/347/497/}{Justia}).

Second, Warren invokes the Court's emerging suspect-classification principle for race. Relying on \textit{Korematsu v. United States}, 323 U.S. 214 (1944), and \textit{Hirabayashi v. United States}, 320 U.S. 81 (1943)---both Fifth Amendment cases involving federal wartime restrictions on Japanese Americans---he reiterates that ``[c]lassifications based solely upon race must be scrutinized with particular care, since they are contrary to our traditions, and hence constitutionally suspect'' (\href{https://supreme.justia.com/cases/federal/us/347/497/}{Justia}). He reinforces this with Justice Harlan's declaration in \textit{Gibson v. Mississippi}, 162 U.S. 565, 591 (1896), that the Constitution forbids racial discrimination in civil and political rights by either the ``general government'' or the states, and with \textit{Buchanan v. Warley}, 245 U.S. 60 (1917), which invalidated a racial residential zoning ordinance as a deprivation of due process (\href{https://supreme.justia.com/cases/federal/us/347/497/}{Justia}).

Third, Warren anchors the holding in substantive due process ``liberty.'' Rejecting a narrow view of liberty as ``mere freedom from bodily restraint,'' he defines ``[l]iberty under law'' to extend to ``the full range of conduct which the individual is free to pursue,'' subject to restriction only for ``a proper governmental objective'' (\href{https://supreme.justia.com/cases/federal/us/347/497/}{Justia}). Because racial segregation in public education is ``not reasonably related to any proper governmental objective,'' it inflicts an arbitrary deprivation of liberty upon Black children in the District of Columbia. Notably, Warren does not repeat \textit{Brown}'s social-science findings about psychological harm, nor does he apply modern strict scrutiny's ``compelling interest'' formulation; he treats school segregation as failing even a baseline requirement of a proper governmental objective.

Fourth, Warren closes with a candid structural and normative appeal: given the Court's holding that same day in \textit{Brown} prohibiting states from maintaining racially segregated public schools, ``it would be unthinkable that the same Constitution would impose a lesser duty on the Federal Government'' (\href{https://www.courtlistener.com/opinion/105222/bolling-v-sharpe/}{CourtListener}; \href{https://supreme.justia.com/cases/federal/us/347/497/}{Justia}). The unanimity of all nine Justices underscored the Court's determination to speak with a single institutional voice on school desegregation nationwide.

\section{Examples: Future Applications}

\begin{enumerate}
  \item \textbf{Same-Side (Invalidated): Federal Sex-Based Social Security Survivor Benefits.} Congress enacts a federal Social Security provision paying survivor benefits to widows with minor children but denying identical benefits to widowers in the exact same position. Under \textit{Bolling}'s Fifth Amendment equal protection component, the federal government cannot enforce sex-based classifications that would violate the Fourteenth Amendment if enacted by a state; because the sex-based distinction is not substantially related to an important governmental objective, it violates Fifth Amendment due process (mirroring \textit{Weinberger v. Wiesenfeld}, 420 U.S. 636 (1975)).

  \item \textbf{Same-Side (Invalidated): Race-Based Federal Contracting Preferences Lacking Narrow Tailoring.} A federal agency adopts a highway subcontracting program that grants financial bonuses to prime contractors who hire minority-owned subcontractors, without demonstrating a compelling remedial interest or narrow tailoring. Under \textit{Bolling} and its progeny (\textit{Adarand Constructors, Inc. v. Pena}, 515 U.S. 200 (1995)), federal racial classifications are subject to the same strict scrutiny as state racial classifications, rendering the program unconstitutional under the Fifth Amendment (\href{https://virginialawreview.org/wp-content/uploads/2020/12/1879.pdf}{Virginia Law Review}).

  \item \textbf{Opposite-Side (Upheld): Federal Alienage Restrictions in Medicare Eligibility.} Congress conditions eligibility for a federal supplemental medical insurance program on five years of continuous lawful residence for noncitizens, whereas a state imposing an analogous residency barrier on welfare benefits would face strict scrutiny under the Fourteenth Amendment. Because Congress possesses plenary authority over immigration and naturalization, \textit{Bolling}'s caveat that due process and equal protection are not ``always interchangeable'' permits the federal government to draw alienage classifications under deferential rational-basis review (as in \textit{Mathews v. Diaz}, 426 U.S. 67 (1976)).

  \item \textbf{Opposite-Side (Upheld): Federal Economic Regulation with Routine Line-Drawing.} Congress imposes a higher federal excise tax on large commercial breweries than on small craft breweries to promote market competition. Because the classification involves neither a suspect class nor a fundamental right, it is reasonably related to a legitimate governmental objective and easily satisfies Fifth Amendment due process under \textit{Bolling}.

  \item \textbf{Fence-Sitter (Close Case): Differential Federal Benefits for Residents of U.S. Territories.} Congress excludes residents of Puerto Rico from a federal Supplemental Security Income (SSI) program available to residents of the fifty states and the District of Columbia. A challenger invoking \textit{Bolling} would argue that the federal government arbitrarily denies equal treatment to U.S. citizens based on geography and, indirectly, ethnicity; the government would counter that the Territory Clause and Puerto Rico's distinct federal tax status supply a rational, proper governmental objective for differential treatment. Because \textit{Bolling} expressly declined to make Fifth and Fourteenth Amendment equality ``always interchangeable,'' courts must weigh structural federalism and territorial powers against anti-discrimination norms (a tension later resolved in favor of rational-basis deference in \textit{United States v. Vaello Madero}, 596 U.S. 159 (2022)).
\end{enumerate}

\section{Critique}

\textit{Bolling v. Sharpe} occupies a paradoxical place in American constitutional law: its result is universally celebrated as a moral imperative, yet its reasoning is widely regarded across the ideological spectrum as among the most textually and historically fragile in the Supreme Court's canon (\href{https://virginialawreview.org/wp-content/uploads/2020/12/1879.pdf}{Virginia Law Review}).

\textbf{Originalist and Textualist Critique:} Originalists and textualists point out two glaring doctrinal flaws. First, the Fifth Amendment was ratified in 1791, when the Constitution explicitly accommodated chattel slavery (e.g., the Fugitive Slave Clause and the Three-Fifths Clause); as scholars such as Richard Fallon and Randy Barnett have observed, virtually no one in 1791 understood the Fifth Amendment's Due Process Clause to prohibit federal racial discrimination (\href{https://virginialawreview.org/wp-content/uploads/2020/12/1879.pdf}{Virginia Law Review}; \href{https://en.wikipedia.org/wiki/Bolling_v._Sharpe}{Wikipedia}). Second, reading an equal protection guarantee into the Fifth Amendment's Due Process Clause creates a severe textual redundancy problem: if ``due process of law'' already encompassed equal protection, the Fourteenth Amendment's separate Equal Protection Clause---placed right next to its own Due Process Clause---would be superfluous. Judge Michael W. McConnell has argued that the Court could have reached the same result without distorting the Fifth Amendment by holding that Congress never actually mandated segregated schools in Washington, D.C., rendering the Board of Education's policy ultra vires (\href{https://en.wikipedia.org/wiki/Bolling_v._Sharpe}{Wikipedia}). Other originalist-leaning scholars, notably Ryan C. Williams and Akhil Reed Amar, have attempted to rescue \textit{Bolling}'s outcome on firmer textual ground by locating a federal equality guarantee in the Fourteenth Amendment's Citizenship Clause, which binds both the federal government and the states (\href{https://virginialawreview.org/articles/originalism-and-other-desegregation-decision/}{Virginia Law Review}).

\textbf{Progressive and Structural Critique:} Progressive scholars likewise identify deep tensions in \textit{Bolling}, though from different angles. Many note that Warren's ``unthinkable'' rationale is pure ipse dixit---a candid concession to pragmatic and moral necessity rather than rigorous constitutional argument---which Richard Primus has described as justified ``by the force of sheer normative necessity'' rather than conventional text or history (\href{https://en.wikipedia.org/wiki/Bolling_v._Sharpe}{Wikipedia}; \href{https://repository.law.umich.edu/cgi/viewcontent.cgi?article=1525&context=facarticles}{Michigan Law Repository}). Furthermore, Warren's revival of broad substantive due process ``liberty'' to strike down arbitrary laws sat uneasily with the post-New Deal Court's repudiation of \textit{Lochner v. New York}, 198 U.S. 45 (1905), and even cited \textit{Buchanan v. Warley}, a \textit{Lochner}-era property-rights decision. Finally, progressive critics argue that turning \textit{Bolling}'s flexible fairness principle into rigid ``congruence'' between the Fifth and Fourteenth Amendments in \textit{Adarand Constructors, Inc. v. Pena} ignored Madisonian structural differences: the Reconstruction Framers distrusted state majorities far more than the national legislature and expressly empowered Congress in Section 5 of the Fourteenth Amendment to enact remedial race-conscious measures (\href{https://virginialawreview.org/wp-content/uploads/2020/12/1879.pdf}{Virginia Law Review}).

\section{Key Quotations}

\begin{quote}
The Fifth Amendment, which is applicable in the District of Columbia, does not contain an equal protection clause, as does the Fourteenth Amendment, which applies only to the states. But the concepts of equal protection and due process, both stemming from our American ideal of fairness, are not mutually exclusive. (\href{https://supreme.justia.com/cases/federal/us/347/497/}{347 U.S. at 499})
\end{quote}

\begin{quote}
The ``equal protection of the laws'' is a more explicit safeguard of prohibited unfairness than ``due process of law,'' and therefore we do not imply that the two are always interchangeable phrases. But, as this Court has recognized, discrimination may be so unjustifiable as to be violative of due process. (\href{https://supreme.justia.com/cases/federal/us/347/497/}{347 U.S. at 499})
\end{quote}

\begin{quote}
Classifications based solely upon race must be scrutinized with particular care, since they are contrary to our traditions, and hence constitutionally suspect. (\href{https://supreme.justia.com/cases/federal/us/347/497/}{347 U.S. at 499})
\end{quote}

\begin{quote}
Liberty under law extends to the full range of conduct which the individual is free to pursue, and it cannot be restricted except for a proper governmental objective. Segregation in public education is not reasonably related to any proper governmental objective, and thus it imposes on Negro children of the District of Columbia a burden that constitutes an arbitrary deprivation of their liberty in violation of the Due Process Clause. (\href{https://supreme.justia.com/cases/federal/us/347/497/}{347 U.S. at 499--500})
\end{quote}

\begin{quote}
In view of our decision that the Constitution prohibits the states from maintaining racially segregated public schools, it would be unthinkable that the same Constitution would impose a lesser duty on the Federal Government. We hold that racial segregation in the public schools of the District of Columbia is a denial of the due process of law guaranteed by the Fifth Amendment to the Constitution. (\href{https://www.courtlistener.com/opinion/105222/bolling-v-sharpe/}{347 U.S. at 500})
\end{quote}

\end{document}
